\documentclass[11pt]{article}
\usepackage[margin=1.05in]{geometry}
\usepackage{amsmath,amssymb,amsthm}
\usepackage{booktabs}
\usepackage{microtype}
\usepackage{xcolor}
\usepackage[colorlinks=true,linkcolor=blue!55!black,citecolor=blue!55!black,urlcolor=blue!55!black]{hyperref}
\usepackage{mathtools}
\usepackage{enumitem}
\setlist{topsep=3pt,itemsep=2pt}

\newcommand{\eps}{\varepsilon}
\newcommand{\dd}{\mathrm{d}}
\newcommand{\Pexp}{\mathbb{P}\exp}
\newcommand{\Q}{\mathbb{Q}}
\newcommand{\Z}{\mathbb{Z}}
\newcommand{\Fp}{\mathbb{F}_p}
\newtheorem{definition}{Definition}[section]
\newtheorem{example}[definition]{Example}
\newtheorem{remark}[definition]{Remark}

\title{Differential Equations for Master Integrals\\[0.4ex]
\large Lecture notes on the analytic evaluation of the double-real
phase-space integrals in $pp \to h + X$ at NNLO}
\date{August 24, 2026}

\begin{document}
\maketitle

\begin{abstract}
\noindent
Integration-by-parts reduction expresses all phase-space integrals of
the double-real channel through a finite basis of master integrals.
These notes develop, from first principles, the method by which the
masters are evaluated: the derivation of first-order differential
equations in the kinematic invariants; the organization of the
equations into triangular systems, and the equivalence relation under
which the many systems collapse to a short list of genuinely distinct
ones; the canonical form of the equations, in which the Laurent
expansion in the dimensional regulator integrates order by order into
iterated integrals; the treatment of square roots of kinematic
polynomials by rational changes of variables and, where none exists,
by the Galois theory of the roots; the determination of boundary
conditions from the analytic and geometric properties of the phase
space; and the recovery of the distributional behaviour of the
integrals at the edges of phase space.  Each concept is defined
before it is used, each mechanism is derived, and each is illustrated
by an example small enough to verify by hand.
\end{abstract}

\tableofcontents

%======================================================================
\section{Introduction}
\label{sec:setting}

\subsection{Phase-space integrals as cut loop integrals}

The physical quantity behind these notes is the NNLO hard function
for single-inclusive hadron production, $pp \to h + X$, in collinear
factorization.  In its double-real channel two unresolved massless
partons are radiated and integrated over, so the basic objects are
phase-space integrals of squared tree amplitudes.  It is convenient
to convert them into loop integrals at the outset.  By the
Sokhotski--Plemelj identity,
\begin{equation}
  \delta^{(+)}(p^2) \;=\;
  \frac{1}{2\pi i}\left[\frac{1}{p^2 - i0} - \frac{1}{p^2 + i0}\right]
  \quad\text{on positive-energy configurations},
\end{equation}
each on-shell delta function is the discontinuity of a propagator;
the positive-energy condition $\theta(p^0)$ is supplied by the cut
itself rather than by the discontinuity.  Every phase-space integral
thereby becomes a two-loop integral in which certain lines are
\emph{cut}.  The value of the rewriting --- this is the device known
as reverse unitarity --- is that cut integrals satisfy the same
integration-by-parts identities as ordinary loop integrals, with one
additional rule: an integral in which a cut propagator has been
cancelled vanishes.  All of loop technology thus becomes available to
phase-space integration.

After projection onto the twist-two correlators, the integrals
depend on the two dimensionless Mandelstam ratios
\begin{equation}
  v = -\frac{t}{s}, \qquad w = -\frac{u}{s},
\end{equation}
whose physical region is the open triangle
\begin{equation}
  0 < v, \qquad 0 < w, \qquad v + w < 1 .
  \label{eq:region}
\end{equation}
The Born configuration sits on the line $v + w = 1$, so
$1 - v - w \to 0$ is the soft limit, while $v \to 0$ and $w \to 0$
are the two collinear limits; all three edges of the triangle carry
physics, and they will reappear throughout, from the singularities of
the differential equations to the boundary conditions and the
endpoint distributions.  Everything is dimensionally regularized,
$d = 4 - 2\eps$.

\subsection{What is to be computed}

Reduction by integration-by-parts identities leaves a basis of $347$
\emph{master integrals} $I_m(v,w;\eps)$.  The reduction is exact
linear algebra over the field $\Q(v,w,\eps)$, but it says nothing
about the values of the basis elements.  The task is their
\emph{analytic evaluation}: for each master, a Laurent expansion
\begin{equation}
  I(v,w;\eps) \;=\; \sum_{k \ge k_0} \eps^{k}\, I^{(k)}(v,w)
\end{equation}
whose coefficients are explicit transcendental functions --- in the
end, $\Q$-linear combinations of iterated integrals of logarithmic
differential forms --- together with exactly known constants.

The method rests on one mathematical fact: \emph{the masters satisfy
first-order linear differential equations in $v$ and $w$}.
Everything else is the systematic exploitation of that fact, and the
chapters build on one another accordingly.
Section~\ref{sec:de} derives the differential equations and their two
structural properties, triangularity and integrability.
Section~\ref{sec:classes} introduces the equivalence relation under
which the many systems of the problem reduce to a short list of
genuinely distinct ones.  Section~\ref{sec:epsform} presents the
canonical form of the equations, shows that it reduces integration to
quadratures, and reads off the local behaviour of the solutions at
the singular loci.  Section~\ref{sec:finding} constructs the
transformation into canonical form.  Section~\ref{sec:roots} treats
the square roots that the kinematics forces on some of the systems.
Section~\ref{sec:paths} assembles the general solution, and
Sections~\ref{sec:boundary} and \ref{sec:endpoint} determine the
boundary data and the behaviour at the edges of phase space.  A
closing section reflects on the respective roles of search and proof,
and a glossary collects the vocabulary.

%======================================================================
\section{Differential equations for master integrals}
\label{sec:de}

\subsection{Masters and integration by parts}

Fix a set of inverse propagators $D_1,\dots,D_N$ --- including the
cut ones --- large enough that every scalar product of loop and
external momenta is a linear combination of the $D_i$ and external
invariants.  Such a set, an \emph{integral family}, defines the
integrals
\begin{equation}
  I_{a_1\dots a_N}(v,w;\eps)
  \;=\; \int \prod_{\ell=1}^{2}\frac{\dd^d k_\ell}{i\pi^{d/2}}\;
  \frac{1}{D_1^{a_1}\cdots D_N^{a_N}},
  \qquad a_i \in \Z ,
  \label{eq:familyint}
\end{equation}
where for a cut line the factor $1/D_i$ stands for the corresponding
delta function, held at unit power throughout; an integral in which a
cut power has been lowered to zero vanishes by the rule of the
introduction.  The \emph{sector} of an integral is the subset of
propagators present with positive power; deleting propagators leads
to \emph{subsectors}, and the sectors of a family form a lattice
under this deletion.  Integration-by-parts identities,
\begin{equation}
  0 \;=\; \int \prod_\ell \dd^d k_\ell\;
  \frac{\partial}{\partial k_j^\mu}
  \left[\frac{q^\mu}{D_1^{a_1}\cdots D_N^{a_N}}\right],
  \qquad q \in \{k_1,k_2,p_i\},
\end{equation}
relate integrals with shifted indices; solving them expresses every
integral of the family through a finite basis of master integrals,
distributed over the sectors of the family.

\subsection{Derivation of the differential system}

Two observations combine into differential equations for the
masters.  A derivative with respect to an invariant can be realized
as a combination of derivatives with respect to external momenta
acting under the integral sign; these raise and lower propagator
powers, producing integrals \emph{of the same family}.  Integration
by parts then maps those integrals back to the masters.  Collecting
the masters of a family and its subsectors into a vector $\vec I$,
one obtains a closed linear system
\begin{equation}
  \partial_v \vec I = A_v(v,w;\eps)\,\vec I,
  \qquad
  \partial_w \vec I = A_w(v,w;\eps)\,\vec I,
  \label{eq:desys}
\end{equation}
with matrices of rational functions.  It is convenient to combine
the pair into the matrix-valued one-form
$A = A_v\,\dd v + A_w\,\dd w$, the \emph{connection}, and to write
$\dd \vec I = A\,\vec I$.

\begin{example}[The one-mass bubble]
\label{ex:bubble}
The mechanism in miniature.  Let
\begin{equation}
  B(s;\eps) = \int \frac{\dd^d k}{i\pi^{d/2}}
  \frac{1}{k^2\,(k+p)^2}, \qquad s = p^2 ,
\end{equation}
a one-scale family with $D_1 = k^2$, $D_2 = (k+p)^2$ and a single
nontrivial master, $B$ itself.  Differentiate with respect to $s$
using $\partial_s = \frac{1}{2s}\, p^\mu \partial_{p^\mu}$:
\begin{equation}
  \frac{\partial B}{\partial s}
  = -\frac{1}{2s}\int \frac{\dd^d k}{i\pi^{d/2}}
    \frac{2p\cdot(k+p)}{k^2\,[(k+p)^2]^2} .
\end{equation}
The numerator is a combination of inverse propagators,
$2 p\cdot(k+p) = D_2 - D_1 + s$, so the right-hand side involves
only integrals of the family: $I_{1,1} = B$, the scaleless
$I_{0,2} = 0$, and $I_{1,2}$.  The integration-by-parts identity
generated by $\partial_{k^\mu}\!\big(k^\mu/D_1 D_2\big)$ reads
$0 = (d-3)\,I_{1,1} + s\,I_{1,2}$, that is,
\begin{equation}
  I_{1,2} = -\frac{d-3}{s}\, I_{1,1},
\end{equation}
and therefore
\begin{equation}
  \frac{\partial B}{\partial s}
  = -\frac{1}{2s}\Big[\,B - (d-3)\,B\,\Big]
  = -\,\frac{4-d}{2s}\,B
  = -\,\frac{\eps}{s}\, B .
  \label{eq:bubblede}
\end{equation}
The equation determines the entire $s$-dependence,
\begin{equation}
  B(s;\eps) \;=\; c(\eps)\,(-s)^{-\eps},
\end{equation}
with the branch fixed by the causal prescription $s \to s + i0$ ---
consistent with dimensional analysis, since $B$ carries mass
dimension $d - 4 = -2\eps$.  A single direct evaluation fixes the
remaining constant,
\begin{equation}
  c(\eps) \;=\;
  \frac{\Gamma(1+\eps)\,\Gamma(1-\eps)^2}
       {\eps\,(1-2\eps)\,\Gamma(1-2\eps)} .
\end{equation}
A differential equation that fixes the functional dependence, and one
boundary constant fixed by an elementary integral: the entire
subject in one line.
\end{example}

\subsection{Triangular structure}

Order the sectors of a family by their number of propagators.
Differentiation and reduction never lead from a sector to one with
more propagators; they lead sideways and down the sector lattice.
With the master vector ordered accordingly, the connection is
\emph{block lower triangular},
\begin{equation}
  A \;=\;
  \begin{pmatrix}
    a_{11} & & & \\
    a_{21} & a_{22} & & \\
    \vdots & & \ddots & \\
    a_{b1} & a_{b2} & \cdots & a_{bb}
  \end{pmatrix},
  \label{eq:blocktriangular}
\end{equation}
where each diagonal entry $a_{ss}$ --- a \emph{block}, in the
terminology fixed in the next section --- couples the masters of one
group of sectors among themselves, and the entries below the
diagonal express how lower sectors source higher ones.  The diagonal
blocks carry the irreducible analytic content of the system; the
entries below the diagonal can be dealt with afterwards --- the
observation on which Section~\ref{sec:finding} rests.

\subsection{Integrability}

Equality of mixed partial derivatives,
$\partial_v\partial_w \vec I = \partial_w\partial_v \vec I$, forces
the compatibility condition
\begin{equation}
  \partial_v A_w - \partial_w A_v - [A_v, A_w] = 0,
  \qquad\text{equivalently}\qquad
  \dd A \;=\; A \wedge A ,
  \label{eq:flatness}
\end{equation}
the statement that the connection is \emph{flat}.  (To check the
sign: $\partial_v\partial_w \vec I = \partial_v(A_w \vec I) =
(\partial_v A_w)\vec I + A_w A_v \vec I$, and antisymmetrizing in
$v,w$ gives $\partial_v A_w - \partial_w A_v + [A_w, A_v] = 0$.)
Flatness plays three roles.  It is a stringent consistency condition
on any derived or transformed system, since it interlocks all matrix
entries.  It guarantees that solutions continued along a path depend
only on the endpoint and on the homotopy class of the path in the
complement of the singular loci, which underlies
Section~\ref{sec:paths}.  And it is stable under change of basis, so
once established for a system it holds for every equivalent form of
that system.

%======================================================================
\section{Families, blocks, and equivalence}
\label{sec:classes}

The double-real channel does not present one large differential
system but many small ones, and many of the small ones turn out to
be the same system in disguise.  Two equivalence relations --- one
among integral families, one among diagonal blocks --- identify the
disguises, and together they set the true size of the analytic
problem.

\subsection{Equivalence of integral families}

Families are generated mechanically, one propagator set per diagram,
and diagrams proliferate: the channel begins with $1898$ of them.
Most are redundant, because different routings of the loop momenta
describe the same set of denominators.

\begin{definition}[Equivalent families; representative families]
Two integral families are \emph{equivalent} if a unimodular integer
redefinition of the loop momenta, composed with shifts by external
momenta, maps one propagator set bijectively onto the other and maps
cut lines to cut lines.  Such a substitution leaves the integration
measure invariant, so equivalent families have identical integrals
up to a relabeling of indices.  A chosen representative of each
equivalence class is called a \emph{representative family}.
\end{definition}

This is where the ``shifting of momenta'' enters the subject.
Equivalence is decided without searching over shift matrices: the
Symanzik polynomials of a family, brought to a normal ordering of
the Feynman parameters by Pak's algorithm (with cut lines
distinguished by auxiliary labels), identify a family up to
relabeling of the parameters.  Coincidence of the normal orderings
is therefore a reliable \emph{candidate} for a momentum shift; every
coincidence is then confirmed by constructing the shift explicitly
and applying it, which is what turns the identification into a
proof.  The $1898$ raw families collapse by better than a factor of
four, and only a fraction of the survivors supports master integrals
at all: ninety-one families, in the end, carry the differential
systems of the problem.

Differentiation is family-local --- the derivative of a master of a
family reduces to masters of that family and its subsectors, never
to unrelated ones --- so the differential-equation problem separates
into $91$ independent systems, of dimensions between $2$ and $47$,
accounting for all $347$ masters.  (The dimensions sum to more than
$347$: each system contains, besides its own masters, the subsector
masters it couples to, and those are counted once, in the family
that owns them.)  The boundaries of the systems were thus fixed
already at the reduction stage; what must be constructed anew is,
for each family, the connection $(A_v, A_w)$, together with the
proof of its flatness.

\subsection{Blocks and their equivalence}

\begin{definition}[Block]
A \emph{block} of a family's system is a strongly connected
component of the coupling structure of its connection: a minimal set
of masters each of which appears in the derivative of another.  With
the masters in sector order, the blocks are the diagonal entries
$a_{ss}$ of Eq.~\eqref{eq:blocktriangular}.  The $91$ family systems
contain $1119$ blocks.
\end{definition}

\begin{definition}[Equivalence of blocks; classes]
\label{def:class}
Two blocks, in the same family or in different ones, are
\emph{equivalent} if their connections coincide as rational matrices
after a permutation of the basis, optionally composed with the
involution $v \leftrightarrow w$.  An equivalence class of blocks is
called a \emph{class}; each class is represented by a single block,
whose canonical form is constructed once and serves for all members.
The $1119$ blocks fall into $173$ classes.
\end{definition}

The roughly tenfold collapse has a physical origin.  A block is the
differential system of the masters of one subtopology, and the same
subtopology occurs as a subdiagram of many different parent
diagrams: the identification of \emph{complete} families by momentum
shifts leaves untouched the far more frequent coincidence of their
\emph{subsectors}.  The involution $v \leftrightarrow w$ is the
crossing symmetry $t \leftrightarrow u$ of the underlying process,
which maps cut diagrams into cut diagrams.  Equivalence is never
assumed on such grounds, however: a candidate identification is
accepted only when the permutation and swap reproduce the
representative's matrices identically.

The consequence for the whole method is that the analysis of
Sections~\ref{sec:finding} and \ref{sec:roots} is carried out once
per class rather than once per block --- a reduction of roughly
tenfold --- after which each family is reconstructed from the
canonical forms of its classes.  The dictionary --- which rows of a
family realize which class, and under which permutation and swap ---
is part of the statement being proved, and is checked again when the
family is put back together.

%======================================================================
\section{The canonical form}
\label{sec:epsform}

\subsection{The $\eps$-form and its immediate consequences}

\begin{definition}[Canonical form; letters, alphabet, residues]
A differential system is in \emph{canonical form} (equivalently, in
\emph{$\eps$-form}; both names are standard) if its connection is
\begin{equation}
  \dd \vec F \;=\; \eps \sum_{a} R_a \,\dd\log \varphi_a(v,w)\;
  \vec F ,
  \label{eq:epsform}
\end{equation}
where the \emph{residue matrices} $R_a$ are constant --- independent
of kinematics and regulator alike --- and the \emph{letters}
$\varphi_a$ are functions of the kinematics alone.  The collection
$\{\varphi_a\}$ is the \emph{alphabet} of the system.  In the
simplest and most common case the letters are polynomials;
Section~\ref{sec:roots} enlarges the notion to algebraic functions.
\end{definition}

One family of the double-real channel, a system of dimension
fourteen, carries the six-letter alphabet
\begin{equation}
  \{\, v,\; w,\; 1-v,\; 1-w,\; v+w,\; 1-v-w \,\} ,
\end{equation}
and every entry of its canonical connection is $\eps$ times a
rational combination of logarithmic differentials of these six
polynomials.  Richer systems have richer alphabets --- among the
two-root families one finds letters quadratic in the variables ---
but always of the same form.

Three consequences follow at once, and they are linked.  The
connection has simple poles exactly on the loci $\varphi_a = 0$ and
nowhere else --- the system is \emph{Fuchsian}, and for generic
$\eps$ each letter with non-vanishing residue is a genuine branch
locus of the solutions, so the alphabet is a complete list of the
places where anything can happen.  Because each
$\dd\log\varphi_a$ is a closed form, the flatness condition
\eqref{eq:flatness} collapses to commutator relations
$\sum_{a<b}\,[R_a,R_b]\;\dd\log\varphi_a \wedge \dd\log\varphi_b =
0$ among constant matrices, evaluated once and for all.  And, as the
next subsection shows, the coefficient of $\eps^k$ in the solution
is a $k$-fold iterated integral, so the expansion is homogeneous in
\emph{transcendental weight} --- the grading on functions and
numbers under which a logarithm carries weight one, a $k$-fold
iterated integral of logarithmic forms weight $k$, and the constants
$\zeta_k$ and $\pi^k$ weight $k$, with rational numbers at weight
zero.  In the canonical basis this grading is manifest, order by
order in $\eps$.

\subsection{Order-by-order integration}

Why one wants the form \eqref{eq:epsform} is best said first.  In a
generic basis the regulator is entangled with the kinematics in
every entry of $A$, and the coefficients of the $\eps$-expansion
obey coupled differential equations with no distinguished structure.
In the canonical basis the regulator counts the number of
integrations and does nothing else: the equations integrate
themselves.

To see this, write $\hat A = \sum_a R_a\,\dd\log\varphi_a$, so that
$\dd\vec F = \eps\,\hat A\,\vec F$ with $\hat A$ independent of
$\eps$.  Substituting $\vec F = \sum_{k}\eps^k \vec F^{(k)}$ and
comparing powers of $\eps$:
\begin{equation}
  \dd \vec F^{(0)} = 0,
  \qquad
  \dd \vec F^{(k)} = \hat A\, \vec F^{(k-1)} \qquad (k \ge 1) .
  \label{eq:orderbyorder}
\end{equation}
The system has disentangled into a recursion of quadratures: the
lowest order is a constant vector, and each subsequent order is an
integral of known functions.  Along a path $\gamma$ from a base
point $p_0$,
\begin{equation}
  \vec F^{(k)}(p) = \vec F^{(k)}(p_0)
  + \int_\gamma \hat A\;\vec F^{(k-1)} ,
\end{equation}
and unwinding the recursion expresses the full solution as the
path-ordered exponential
\begin{equation}
  \vec F(p) \;=\; \Pexp\!\Big(\eps \int_{\gamma} \hat A\Big)\,
  \vec F(p_0)
  \;=\; \sum_{k\ge0} \eps^k \!\!
  \sum_{a_1,\dots,a_k}\!\!
  R_{a_k}\!\cdots R_{a_1}
  \int\limits_{0\le t_1\le\cdots\le t_k\le 1}
  \prod_{i=1}^{k}\dd\log \varphi_{a_i}\big(\gamma(t_i)\big)\;
  \vec F(p_0) .
  \label{eq:pexp}
\end{equation}
The $k$-fold integrals are \emph{Chen iterated integrals} over the
alphabet; flatness makes them independent of the choice of path
within its homotopy class.  The analytic problem is thereby solved
in principle, and what remains is to express the iterated integrals
in a standard basis of functions (Section~\ref{sec:paths}) and to
determine the boundary vector $\vec F(p_0)$
(Section~\ref{sec:boundary}).

One apparent paradox is worth resolving immediately.  The masters
themselves have poles in $\eps$ --- their Laurent expansions begin
at negative orders --- while the canonical recursion
\eqref{eq:orderbyorder} starts from a \emph{constant} at order
$\eps^0$.  There is no contradiction: the poles reside in the
transformation between the two bases.  The matrix $T$ relating the
derived masters to the canonical ones carries inverse powers of
$\eps$ from the normalization of the basis, so a finite canonical
order feeds several Laurent orders of a master; statements about
uniform weight refer to the canonical basis.

\begin{example}[A one-dimensional system]
$\dd F = \eps\,\dd\log w\,F$ has the solution
$F = C\,w^{\eps} = C\big(1 + \eps\log w + \tfrac{\eps^2}{2}\log^2 w
+ \cdots\big)$: the recursion \eqref{eq:orderbyorder} generates
exactly the expansion of the power function.  Both descriptions of
this object --- the expanded logarithmic series and the unexpanded
power $w^\eps$ --- will be needed; the tension between them is the
subject of Section~\ref{sec:endpoint}.
\end{example}

\begin{example}[A two-dimensional system and the dilogarithm]
\label{ex:li2}
Take the alphabet $\{w,\,1-w\}$ on $0<w<1$ and the system
\begin{equation}
  \dd \vec F = \eps\left[
  \begin{pmatrix} 0 & 0\\ 0 & -1\end{pmatrix}\dd\log w
  +
  \begin{pmatrix} 0 & 0\\ 1 & 0 \end{pmatrix}\dd\log(1-w)
  \right]\vec F ,
\end{equation}
with the boundary condition $\vec F \to (1,0)^{\mathsf T}$ as
$w \to 0$ (the precise sense in which a value is assigned at a
singular point is given in Section~\ref{sec:local}).  The first
component stays constant.  For the second, the recursion gives
\begin{align}
  F_2^{(1)}(w) &= \int_0^w \dd\log(1-w')\cdot 1
               \;=\; \log(1-w) ,\\
  F_2^{(2)}(w) &= \int_0^w \Big[\dd\log(1-w')\cdot 0
               \;-\;\dd\log w' \cdot F_2^{(1)}(w')\Big]
               = -\int_0^w \frac{\dd w'}{w'}\,\log(1-w')
               \;=\; \operatorname{Li}_2(w) .
\end{align}
The dilogarithm appears as the twofold iterated integral with the
word $(w,\,1{-}w)$; the next order produces trilogarithms, and so
on.  Every master of the calculation is, order by order, a rational
combination of such words in its own alphabet.
\end{example}

\subsection{Local behaviour at the singular loci}
\label{sec:local}

The canonical form does more than integrate: it displays, before any
integration, how every solution behaves at every singular locus.
Near a zero of a letter --- say $w \to 0$, with letter $w$ --- the
system reads
\begin{equation}
  \partial_w \vec F \;=\; \frac{\eps\,R_{w}}{w}\,\vec F
  \;+\; (\text{regular}) ,
\end{equation}
and the classical Frobenius analysis applies: for generic $\eps$ a
fundamental system of local solutions is
\begin{equation}
  \vec F \;\sim\; w^{\,\eps R_w}\,\vec h(w)
  \;=\; \sum_{\text{modes}} w^{\,n + a\eps}\,\big(\log w\big)^{j}\,
  \big(\text{const} + O(w)\big) ,
  \label{eq:frobenius}
\end{equation}
with the coefficients $a$ ranging over the eigenvalues of $R_w$,
non-negative integer shifts $n$ from the Taylor expansion of the
regular part, and powers of $\log w$ from the Jordan structure of
$R_w$.  \emph{The complete local behaviour of every solution at
every singular locus is thus encoded in the residue matrix at the
corresponding letter} --- exponents in its eigenvalues, logarithms
in its Jordan blocks --- and is known exactly in advance.  Matching
the global solution onto \eqref{eq:frobenius} is also what defines a
boundary value at a singular point: the constant in the mode of
exponent zero.

Two qualifications keep the statement honest.  The normal form
requires that no two exponents $\eps\lambda_i$, $\eps\lambda_j$
differ by a nonzero integer --- automatic for generic $\eps$, which
is the regime of the Laurent expansion, but not an identity in
$\eps$.  And ``logarithms come from the Jordan blocks'' refers to
the solution at fixed $\eps$; upon expanding in $\eps$ the
logarithms proliferate for a second, independent reason,
\begin{equation}
  w^{\,\eps R_w} \;=\; e^{\,\eps R_w \log w}
  \;=\; \sum_{m \ge 0} \frac{(\eps \log w)^m}{m!}\, R_w^{\,m} ,
\end{equation}
so a power of $\log w$ appears at every order in $\eps$ whether or
not $R_w$ is diagonalizable; the Jordan structure bounds only the
logarithms that survive at fixed $\eps$.  This distinction --- the
unexpanded mode versus its logarithmic series --- is precisely what
the endpoint analysis of Section~\ref{sec:endpoint} turns on.

\begin{remark}[Existence]
No theorem guarantees that a given block admits a canonical form
with rational letters.  In the double-real channel every one of the
classes admits one --- though for some only after a change of
variables that rationalizes a square root, and for the classes tied
to three independent roots the appropriate generalization is the
subject of Section~\ref{sec:roots}.  The methodological position
throughout is that the canonical form is a \emph{statement to be
verified}: any means whatsoever may produce the transformation, and
what establishes it is the exact identity between the transformed
connection and the form \eqref{eq:epsform}.
\end{remark}

%======================================================================
\section{Construction of the canonical transformation}
\label{sec:finding}

The unknown of this section is an invertible matrix $T(v,w;\eps)$
relating the derived basis to the canonical one,
$\vec I = T\,\vec F$.  Substituting into $\dd\vec I = A\vec I$ shows
that the connection transforms inhomogeneously,
\begin{equation}
  A \;\longmapsto\; A' \;=\; T^{-1} A\, T \;-\; T^{-1}\dd T ,
  \label{eq:gauge}
\end{equation}
and the requirement is that $A'$ be of the form \eqref{eq:epsform}.
The triangular structure \eqref{eq:blocktriangular} decides the
order of the work.  A diagonal block never feels what lies below it,
so each block can be brought to canonical form on its own --- and
since the same block recurs across many families, this is done once
for each class.  With the diagonal canonical, the entries below it
yield to transformations of a particularly simple, unipotent shape.
What survives both steps is a dependence on $\eps$ that is constant
in the kinematics, and it is removed last, by conjugation.

\subsection{Diagonal blocks: Fuchsian reduction and balances}
\label{sec:diagonal}

For the diagonal blocks two complementary strategies are classical,
and each produces a candidate transformation which then stands or
falls on the exact identity.

The first is the \emph{rational ansatz}: posit for $T$ a matrix of
rational functions whose denominators are drawn from the expected
letters and whose numerators are polynomials of bounded degree with
undetermined $\eps$-dependent coefficients; impose that
\eqref{eq:gauge} produce the form \eqref{eq:epsform}; solve the
resulting linear equations for the coefficients.  The approach is
exhaustive at a given degree.

The second proceeds through normal forms.  One first reduces the
connection to \emph{Fuchsian} form --- simple poles only --- and
then normalizes the eigenvalues of the residue matrices by
\emph{balance transformations} between a pair of singular points
$v_1, v_2$,
\begin{equation}
  T \;=\; \bar P + \frac{v - v_1}{\,v - v_2\,}\, P ,
  \qquad \bar P = 1 - P ,
\end{equation}
with $P$ the spectral projector of the residue at $v_1$ onto the
eigenspace being moved.  Such a transformation lowers the chosen
eigenvalue at $v_1$ by one and raises the matching eigenvalue at
$v_2$ by one, leaving every other singular point untouched; a
balance always acts at two points, never one, since the sum of all
residue eigenvalues over the singular points (infinity included) is
invariant.  A sequence of balances aligns all local exponents
proportionally to $\eps$, after which a constant transformation
factors the regulator out of the connection.

\begin{example}[What a higher pole means, and what is removable]
\label{ex:fuchs}
Consider first the scalar equation
\begin{equation}
  \frac{\dd F}{\dd w} = \left(\frac{\alpha}{w^2} +
  \frac{\beta\,\eps}{w}\right) F .
\end{equation}
The substitution $F = e^{-\alpha/w}\, G$ does remove the double
pole, since $\frac{\dd}{\dd w} e^{-\alpha/w} =
\frac{\alpha}{w^2}\,e^{-\alpha/w}$, leaving
$G' = (\beta\eps/w)\,G$ and $G = C\,w^{\beta\eps}$.  But
$e^{-\alpha/w}$ is not a rational --- not even a meromorphic ---
change of basis, and no rational one can do the job: for a scalar
equation a rational $F = r(w)\,G$ shifts the connection by
$-\dd\log r$, which has only simple poles with integer residues.  A
scalar double pole is therefore always a genuinely irregular
singularity, with the essential-singular factor $e^{-\alpha/w}$ in
its solution, and no canonical form exists.

What \emph{is} removable is a higher pole sitting in a nilpotent
direction of a matrix residue.  For
\begin{equation}
  A \;=\; \frac{1}{w^2}
  \begin{pmatrix} 0 & \alpha \\ 0 & 0 \end{pmatrix} \dd w ,
  \qquad
  T = \begin{pmatrix} 1 & 0 \\ 0 & w \end{pmatrix}
  \quad\Longrightarrow\quad
  A' \;=\; \frac{1}{w}
  \begin{pmatrix} 0 & \alpha \\ 0 & -1 \end{pmatrix} \dd w ,
\end{equation}
a simple pole, reached by a rational transformation --- Fuchsian,
and ready for the balances above.  In the present calculation the
genuinely irregular case never occurs once the variables of
Section~\ref{sec:roots} are chosen correctly.
\end{example}

\subsection{Coupling the blocks: the completion equation}
\label{sec:completion}

Suppose the diagonal blocks are already canonical --- each equal to
$\eps\,e_k$ with $e_k$ a constant-residue logarithmic form ---
while the entries below the diagonal, denoted $b_{kj}$ with $j<k$,
are not.  They are brought to canonical form one block row at a
time, by transformations of unipotent shape:
\begin{equation}
  T = 1 + D, \qquad
  D\ \text{supported on block row } k,\ \text{columns } j<k .
\end{equation}
Because $D$ is strictly triangular, $D^2 = 0$ and $T^{-1} = 1 - D$;
moreover $D A D$ and $D\,\dd D$ vanish --- the latter because $D$
and $\dd D$ share the same support, the former because $A$ is block
lower triangular, so its $(j,k)$ entries with $j < k$ are zero ---
and the transformation law \eqref{eq:gauge} is exact without
expansion:
\begin{equation}
  A' \;=\; A + A D - D A - \dd D .
\end{equation}
Reading off the $(k,j)$ block of the requirement that $A'$ be
canonical gives the \emph{completion equation}
\begin{equation}
  \dd D_{kj}
  \;=\; \eps\big( e_k\, D_{kj} - D_{kj}\, c_j \big)
  \;+\; \bar b_{kj}
  \;-\; \eps \sum_a K_a\, \dd\log\varphi_a ,
  \label{eq:completion}
\end{equation}
in which $c_j$ is the canonical form of the diagonal block in
column $j$ (so that $a_{jj} = \eps\,c_j$, in parallel with
$a_{kk} = \eps\,e_k$), the source
$\bar b_{kj} = b_{kj} - \sum_{j<m<k} D_{km}\, b_{mj}$ collects the
original coupling together with the contributions of blocks of the
same row completed previously, and the constant matrices $K_a$ are
unknowns alongside $D_{kj}$: they become the residues of the
completed entry.  The essential feature of \eqref{eq:completion} is
its \emph{linearity} in the unknowns $(D_{kj}, K_a)$, which places
the whole completion within reach of linear algebra.

\begin{example}[A completion displaying all three mechanisms]
\label{ex:completion}
In one variable, with diagonal data $e = \dd\log v$ and $c = 0$,
take the coupling
\begin{equation}
  b \;=\; \eps\,\dd v \;+\; \frac{\dd v}{1+v} .
\end{equation}
The two pieces are absorbed by different mechanisms.  For the first,
try a rational $D$: with $D = \dfrac{\eps\, v}{1-\eps}$ one has
\begin{equation}
  \eps\, e\, D + \eps\,\dd v
  = \eps\,\frac{\dd v}{v}\cdot\frac{\eps v}{1-\eps} + \eps\,\dd v
  = \frac{\eps}{1-\eps}\,\dd v
  = \dd D ,
\end{equation}
so the non-logarithmic part of the coupling disappears into the
transformation.  The second piece is itself a logarithmic form with
constant coefficient, $\dd\log(1+v)$: no rational $D$ can absorb it
(that would require $\log(1+v)$ itself), and none should --- it is
kept as a residue, $K_{1+v} = 1$.  It appears at order $\eps^0$
rather than multiplied by $\eps$, but that mismatch is a
normalization of the basis, constant in the kinematics, and is
removed by the final transformation of
Section~\ref{sec:constT}.  Rational parts of a coupling are
absorbed by $D$; logarithmic parts with constant coefficients become
residues; normalization mismatches in $\eps$ are deferred to the
end --- this division of labour persists in the general case.
\end{example}

\subsection{Solution of the completion equation by modular
arithmetic}
\label{sec:modular}

For the systems of interest the unknown $D_{kj}$ may have several
hundred rational-function entries of high degree, and solving
\eqref{eq:completion} symbolically is prohibitive.  The equation is
linear, however, and a linear problem can be solved numerically over
finite fields and reassembled exactly.  One writes
$D_{kj} = N(v,w;\eps)\big/\prod_a \varphi_a^{m_a}$, with the
denominator a product of letters and the numerator a polynomial with
undetermined coefficients, the admissible monomials being bounded by
a degree analysis of the source $\bar b_{kj}$.  Substituting
rational values of $(v,w)$ \emph{and of the regulator}, all reduced
modulo a large prime $p$, turns \eqref{eq:completion} into exact
linear equations over the finite field $\Fp$ for the unknown
coefficients; enough substitutions determine them modulo $p$, and
the dependence on $\eps$ --- the solution is a rational function of
the regulator as well --- is recovered by sampling sufficiently many
regulator values and interpolating.  Repeating over several primes,
combining by the Chinese remainder theorem, and lifting each
coefficient to a rational number by rational reconstruction yields
an exact candidate, confirmed on evaluations withheld from the fit
and at a fresh prime.  None of this is a proof; the proof is the
substitution of the reconstructed $(D, K)$ back into
\eqref{eq:completion} --- or into the equivalent statement for the
fully assembled family --- and the exact verification of the
identity.

\begin{example}[Rational reconstruction]
Take $p = 101$ and suppose a coefficient is determined as
$29 \bmod 101$.  Rational reconstruction seeks a fraction of small
numerator and denominator congruent to it: since
$7 \cdot 29 = 203 = 2\cdot 101 + 1 \equiv 1 \pmod{101}$, one has
$29 \equiv \tfrac{1}{7} \pmod{101}$.  A second prime, $p = 103$,
where $\tfrac17 \equiv 59 \pmod{103}$ (as $7 \cdot 59 = 413 =
4\cdot103+1$), and the Chinese remainder lift of $(29, 59)$ modulo
$101\cdot103$ reconstructs $\tfrac17$ again.  The guarantee behind
the search is a lattice bound: modulo a single prime $p$ the
reconstruction is unique only for numerator and denominator below
$\sqrt{p/2}$ --- about $7.1$ for $p = 101$, so $\tfrac17$ is
admissible there only just, which is exactly why the second prime is
not optional.  Stability of the lift as the modulus grows, plus
agreement at a prime never used, is compelling evidence; the proof
remains the final exact substitution.
\end{example}

\begin{remark}[Non-uniqueness of the completion]
\label{rem:representative}
When the linear system for $(D, K)$ has a null space the completed
entry is not unique; all solutions are related by redefinitions that
preserve the canonical form.  All are correct, but they differ
substantially in the degrees of the rational functions they insert
into the sources $\bar b$ of the rows still to be completed, and
degree growth compounds from row to row.  A definite and favourable
choice results from requiring the coefficients of the highest
admissible numerator monomials to vanish, which selects a
representative of low degree.  The choice leaves no trace in the
final canonical form; it decides only the degrees that the remaining
rows inherit.
\end{remark}

\subsection{The residual dependence on the regulator}
\label{sec:constT}

After the completion, the connection is logarithmic with
$\eps$-independent letters, but its residue matrices may still
depend on $\eps$ --- rational functions of $\eps$ in place of the
required constants times $\eps$.  The origin of this residue is the
normalization of the master basis: each sector's masters were
defined with $\eps$-dependent prefactors, and a change of
normalization is a transformation \emph{constant in the
kinematics}.  For constant $T(\eps)$ the inhomogeneous term of
\eqref{eq:gauge} vanishes and the transformation acts by
simultaneous conjugation of all residues:
\begin{equation}
  T^{-1}(\eps)\, \tilde R_a(\eps)\, T(\eps) \;=\; \eps\, R_a
  \qquad \text{for all } a .
\end{equation}
Finding such a $T$ is a finite-dimensional problem: evaluate the
connection at a few exact rational kinematic points, factor the
$\eps$-dependence of the resulting constant matrices, and verify
the candidate by the exact identity on the full connection.
Composing the class transformations, the row completions, and this
constant conjugation yields the matrix $T$ that carries the derived
system \eqref{eq:desys} into the canonical form
\eqref{eq:epsform}.

%======================================================================
\section{Algebraic letters: square roots and their rationalization}
\label{sec:roots}

\subsection{Square roots from the geometry of phase space}

Nothing so far required the letters to be polynomials, and for a
minority of the blocks they are not.  The origin is geometric.
Certain sub-configurations of the three-particle final state
degenerate on the vanishing of a Gram determinant, which in the
variables $(v,w)$ is the K\"all\'en polynomial
\begin{equation}
  \lambda_1(v,w) \;=\; (1 - v - w)^2 - 4\,v\,w
  \;=\; \lambda(1, v, w) ,
  \label{eq:kallen}
\end{equation}
where $\lambda(a,b,c) = a^2+b^2+c^2-2ab-2bc-2ca$ is the triangle
function.  Its zero locus is the curve $\sqrt v + \sqrt w = 1$, a
pseudo-threshold lying \emph{inside} the physical triangle
\eqref{eq:region} rather than on its boundary --- the point
$v = w = \tfrac14$ sits on it --- and it separates the region
$\lambda_1 > 0$, where $\sqrt{\lambda_1}$ is real, from the rest.
Masters of sectors sensitive to this locus have branch points of
square-root type there, and a square-root branch point cannot be
produced by logarithmic differentials of polynomials alone: the
alphabet must contain $\sqrt{\lambda_1}$, typically through letters
of the form
\begin{equation}
  \frac{1 - v - w - \sqrt{\lambda_1}}
       {\,1 - v - w + \sqrt{\lambda_1}\,} ,
\end{equation}
which are exchanged with their reciprocals under the sign change
$\sqrt{\lambda_1} \to -\sqrt{\lambda_1}$ (in the language of
covering spaces, the deck transformation of the double cover).  The
double-real channel realizes six such radicands:
\begin{equation}
  \lambda_1,\qquad
  \lambda_2 = \lambda_1(-v,w),\qquad
  \lambda_3 = \lambda_1(v,-w),\qquad
  4v + w^2,\qquad v^2 + 4w,\qquad 1 - 4vw .
  \label{eq:radicands}
\end{equation}
Under the involution $v \leftrightarrow w$, the pairs
$\{\lambda_2, \lambda_3\}$ and $\{4v{+}w^2,\, v^2{+}4w\}$ are
exchanged, while $\lambda_1$ and $1-4vw$ are invariant.

\begin{definition}[Root content of a system]
Radicands are counted modulo perfect-square factors and modulo
products: $\sqrt{q_1}$, $\sqrt{q_2}$ and $\sqrt{q_1 q_2}$ constitute
\emph{two} independent roots, since the third is determined by the
other two.  Formally, the radicands of a system generate a subgroup
of the square-class group
$\Q(v,w)^\times / \big(\Q(v,w)^\times\big)^2$, and the \emph{root
content} is the rank of that subgroup.  In the double-real channel
most families involve no root at all; a third involve one; a dozen
involve two; and three families --- the hard case treated in
Section~\ref{sec:galois} --- involve three.
\end{definition}

The root content is read off the derived connection and reduced
exactly; it is the single datum that determines which of the
techniques below applies.

\subsection{Rationalization by change of variables}

\begin{definition}[Rationalizing substitution]
A \emph{rationalizing substitution} for a radicand $q(v,w)$ is a
rational change of variables $(v,w) = \big(f(x,y),\,g(x,y)\big)$
with non-degenerate Jacobian such that $q(f,g) = r(x,y)^2$ for a
rational function $r$.  Under such a substitution $\sqrt q$ becomes
the rational function $r$, and a system whose alphabet involves only
$\sqrt q$ acquires polynomial letters in $(x,y)$.
\end{definition}

\begin{example}[Rationalizing the K\"all\'en root]
\label{ex:kallenchart}
Consider the substitution suggested by the factorized form of the
degenerate configurations,
\begin{equation}
  v = x\,y, \qquad w = (1-x)(1-y) .
\end{equation}
Then
\begin{equation}
  1 - v - w = x + y - 2xy,
  \qquad
  \lambda_1 = (x+y-2xy)^2 - 4\,xy\,(1{-}x)(1{-}y)
  = (x-y)^2 ,
\end{equation}
by direct expansion.  Hence $\sqrt{\lambda_1} = x - y$, the
Jacobian determinant is $\partial(v,w)/\partial(x,y) = x - y$,
non-vanishing away from the diagonal, and the algebraic letter above
becomes
\begin{equation}
  \frac{(x+y-2xy)-(x-y)}{(x+y-2xy)+(x-y)}
  \;=\; \frac{y\,(1-x)}{x\,(1-y)} ,
\end{equation}
a ratio of polynomial letters.  Two facts about the substitution
deserve recording.  It is two-to-one: both $v$ and $w$ are symmetric
under $x \leftrightarrow y$, and that exchange is precisely the deck
transformation $\sqrt{\lambda_1} \to -\sqrt{\lambda_1}$, branched on
the diagonal $x = y$ where the Jacobian vanishes.  And its real
image is not the whole triangle \eqref{eq:region} but the
sub-region $\sqrt v + \sqrt w \le 1$, i.e.\ $\lambda_1 \ge 0$;
outside it $\sqrt{\lambda_1}$ is imaginary, and reaching it is an
analytic continuation to be performed deliberately.  Paths chosen in
$(x,y)$ therefore live in that sub-region.  The sign variants
$v = \pm xy$, $w = \mp(1-x)(1-y)$ rationalize $\lambda_2$ and
$\lambda_3$ in the same way, and substitutions of the same conic
type rationalize $4v+w^2$, $v^2+4w$ and $1-4vw$ --- one
substitution for each of the six radicands.
\end{example}

The construction behind all such substitutions is classical: the
equation $q(v,w) = r^2$ defines a conic in suitable projective
coordinates, and a conic possessing one rational point is rationally
parametrized by the pencil of lines through that point --- the same
device that parametrizes the Pythagorean triples.  For each radicand
of \eqref{eq:radicands} the substitution can be written down
explicitly, together with its root image and Jacobian, and the
identity $q\circ(f,g) = r^2$ is then a matter of expansion, never
taken on faith.

Working in the new variables requires one transformation of the
whole system by the chain rule,
\begin{equation}
  A_x = A_v\,\partial_x f + A_w\,\partial_x g,
  \qquad
  A_y = A_v\,\partial_y f + A_w\,\partial_y g ,
\end{equation}
with flatness re-established in $(x,y)$; afterwards the
constructions of Sections~\ref{sec:epsform}--\ref{sec:finding}
apply verbatim.

\subsection{Several roots}

A system whose alphabet involves two independent roots requires a
substitution rationalizing both at once.  Such substitutions are
found iteratively: rationalize the first root as above; under the
substitution, the second radicand becomes a new polynomial
$q_2(x,y)$, defining a second conic \emph{in the new variables}; a
rational point on that conic yields a parametrization by lines, in
one of the variables, that rationalizes it while preserving the
rationality of the first root.  Three iterated substitutions of this
kind --- one for each pair of K\"all\'en-type radicands --- suffice
for every two-root family that occurs, and the resulting canonical
connections are naturally written in the final pair of parameters.

\subsection{The Galois algebra of the roots}
\label{sec:galois}

For the families involving three independent roots, no rational
substitution can rationalize all three simultaneously; this is a
theorem, not a failure of search.  A rationalizing substitution is a
dominant rational map from the plane, so its existence would make
the fibred product of the three conic double covers a unirational
surface; for the root triples occurring here that surface has two
ordinary double points whose resolution is a K3 surface, which is
not unirational in characteristic zero.  The obstruction is
geometric and definitive.

The response is to renounce parametrization and compute in the
extension the roots generate.  Adjoin them as symbols:
\begin{equation}
  E \;=\; \Q(v,w,\eps)\,[\,r_1,r_2,r_3\,]\ \big/\
  \big(\,r_i^2 - q_i\,\big) ,
\end{equation}
an algebra of dimension $2^3$ with basis
$r_S = \prod_{i\in S} r_i$ indexed by subsets
$S \subseteq \{1,2,3\}$; when the $q_i$ are independent square
classes --- which is what root content three means --- $E$ is a
field extension of degree eight.  Its Galois group is $(\Z/2)^3$,
acting by the independent sign changes $r_i \mapsto -r_i$.  Every
element decomposes into the characters of this group (its
\emph{grades}); multiplication respects the grading,
\begin{equation}
  r_S\, r_T \;=\; \Big(\prod_{i \in S\cap T} q_i\Big)\;
  r_{S \mathbin{\triangle} T} ,
\end{equation}
and differentiation acts grade-wise through
$\dd r_S = \tfrac12\, r_S \sum_{i\in S} \dd\log q_i$.  All linear
problems of Section~\ref{sec:finding} --- the completion equation,
its evaluation over finite fields, the reconstruction --- can be
posed on the $8$-component coefficient vectors, whose entries are
ordinary rational functions.  Evaluations modulo a prime are
performed at points where every $q_i$ is a nonzero quadratic
residue, so that all eight embeddings of the algebra into $\Fp$
exist and the grade decomposition becomes a discrete Fourier
transform over $(\Z/2)^3$; agreement of a solution across all eight
sign embeddings is a consistency requirement that a correct
computation must, and does, satisfy.  No better description of these
systems is presently known; the construction has been carried
through on model systems of every root rank, and the physical
three-root families are of the same type.

\begin{remark}[Roots and the crossing involution]
The involution $v \leftrightarrow w$ that enters the equivalence of
blocks (Definition~\ref{def:class}) exchanges $\lambda_2$ with
$\lambda_3$ and $4v+w^2$ with $v^2+4w$ while fixing $\lambda_1$ and
$1-4vw$; within the K\"all\'en substitution it takes the simple form
$(x,y) \mapsto (1-x,\,1-y)$.
\end{remark}

%======================================================================
\section{The general solution: paths and polylogarithms}
\label{sec:paths}

The canonical form reduces integration to the evaluation of the
iterated integrals in \eqref{eq:pexp} along paths in the physical
region.  For letters \emph{linear} in the integration parameter the
iterated integrals belong to a standard class:

\begin{definition}[Generalized polylogarithms]
For $a_1,\dots,a_k \in \mathbb{C}$ define recursively
\begin{equation}
  G(a_1,a_2,\dots,a_k;\,t) \;=\;
  \int_0^{t}\frac{\dd t'}{t' - a_1}\;
  G(a_2,\dots,a_k;\,t') ,
  \qquad G(;t) = 1 ,
\end{equation}
with the logarithmic case
$G(\vec 0_k; t) = \tfrac{1}{k!}\log^k t$.  These \emph{generalized
polylogarithms} of weight $k$ contain the classical functions:
$G(0,1;t) = -\operatorname{Li}_2(t)$, matching
Example~\ref{ex:li2} up to the sign conventions, and all functions
of weight two and three reduce to classical polylogarithms.
\end{definition}

Evaluating a family is then a matter of choosing a base point $p_0$
in the triangle \eqref{eq:region} and a path to the point of
interest, composed of segments on which every letter restricts to a
polynomial of degree one in the path parameter --- for then the
solution along each segment is a combination of the functions $G$.
When a letter restricts to a quadratic, subdividing the segment does
not help: the restriction stays quadratic on every piece.  The
remedy is either to choose the segment differently --- in practice
segments parallel to the coordinate axes, on which the bilinear
letters produced by the rationalizing substitutions become linear
again --- or to factor the quadratic and admit its two linear
factors, with algebraic constants, as letters in their own right.
Flatness guarantees that the result depends only on the endpoint and
on the homotopy class of the path; the zero locus of a letter is
crossed only as a deliberate analytic continuation, never in
passing.  Families whose canonical form lives in rationalizing
variables are evaluated in those variables throughout, along paths
chosen by the same linearity criterion, and only the final values
are mapped back to $(v,w)$.

\subsection{The order of expansion required}

The cost of an iterated integral grows steeply with its weight, so
one does not expand every master to the same depth; the required
depth is dictated by where the master enters the cross section.  A
master multiplying a double pole of a coefficient needs two orders
in $\eps$ beyond one multiplying a finite coefficient, and the
required order is found for each master by following its
coefficients through the assembled combination.  In the double-real
channel the outcome is lopsided: of the $347$ masters, $203$ are
needed only to their finite order and $131$ one order deeper, while
$9$ are needed deeper still and the remaining four less deeply (two
enter only at order $\eps$, and two multiply coefficients that
vanish identically).  This accounting is settled before any
integration, and it sets the scale of the entire computation.

%======================================================================
\section{Boundary conditions}
\label{sec:boundary}

\subsection{The space of solutions and the conditions upon it}

The path-ordered solution \eqref{eq:pexp} involves the undetermined
vector $\vec F(p_0) = \sum_k \eps^k\,\vec C^{(k)}$: one constant
per master per order.  Determining these constants is where analysis
meets physics, and the governing principle is that almost all of
them are fixed by consistency, without evaluating a single new
integral.

The most productive condition --- derived in the next subsection ---
is the \emph{exclusion of local behaviour}.  Near each edge of the
physical triangle the differential equation offers local modes with
exponents read from its residue matrices
(Section~\ref{sec:local}); the integral representation of the
master permits only some of them; every offered mode that the
representation forbids annihilates one combination of constants.

Three further conditions complete the analysis.  Elementary
normalizations: the simplest master of the lowest sector is the
phase-space volume itself, computable in closed form
(Example~\ref{ex:volume}), and its expansion provides initial data
wherever that sector appears.  Inheritance: families are analysed in
an order compatible with the subsector relation, and a master of a
subsector, once known, is known in every family in which that
subsector reappears --- the identification of
Section~\ref{sec:classes} is an equality of integrals, not an
analogy --- imposing linear conditions there.  Finiteness at
spurious loci: some zero loci of letters lie inside the region where
a given integral is manifestly finite, and the local modes singular
there must then be absent from it.

\subsection{Power counting at the boundary}
\label{sec:powercounting}

Consider a master given by a convergent cut phase-space integral,
and examine one edge of the triangle, say $w \to 0$.  In the limit
the integral is governed by a small set of momentum regions ---
hard, soft, collinear to one or the other direction --- and each
region contributes to the integral with a definite overall scaling
in $w$.  Consequently the exact function has an expansion in the
modes
\begin{equation}
  w^{\,n + a\eps},
  \qquad n \in \Z_{\ge n_0} ,
\end{equation}
where the admissible values of $a$ form a short list fixed by the
region analysis: which lines of the sector can become soft or
collinear in the limit, and with what measure.  It is a coarse datum
--- a handful of integers and multiples of $\eps$, common to whole
groups of sectors --- though it does depend on the propagator
content of the sector, since that is what decides which regions
exist.  (Over the whole double-real channel, for instance, the
spectrum of $\eps$-coefficients on the soft edge consists of just
five values.)  The differential equation, on the other hand, offers
local modes with exponents drawn from the eigenvalues of the residue
matrix at the corresponding letter.  Whenever the equation offers an
exponent that the power counting excludes, the coefficient of the
entire corresponding mode must vanish: one exact, homogeneous,
linear condition on the boundary constants per excluded mode.  The
same analysis applies at $v \to 0$ and on the soft edge
$1 - v - w \to 0$, with the distance to the edge as the scaling
variable.  Two cautions are essential in practice.  The admissible
exponents must be derived in the same variables in which the local
analysis is performed --- a rationalizing substitution alters the
measure by its Jacobian and shifts the integer parts --- and
degenerate (Jordan) directions of the residue matrix couple modes
logarithmically, so the exclusion argument must be applied to
generalized eigenspaces, not to eigenvalues alone.

\subsection{The kernel of the constraint system}

All conditions on a given block assemble into a linear system on
its boundary constants; call its matrix $C$.  The number of
genuinely undetermined constants is $\dim \ker C$, and it is known
before a single integral is evaluated: it is a property of the
constraint matrix, not something that emerges as the calculation
proceeds.  Constants surviving in the kernel are then identified
across blocks --- the same subsector appearing in several families
contributes one constant, not several --- which reduces the count
over the whole problem to a short list of genuinely new numbers.

\begin{definition}[Boundary periods]
The surviving constants are values of cut phase-space integrals at
distinguished limits of the kinematics, taken in a definite order:
\emph{boundary periods}.  At the weights relevant here they evaluate
to rational combinations of $\zeta$-values and kindred constants.
Each is established by an exact derivation --- through a
Mellin--Barnes representation (trading the integrand for a contour
integral of $\Gamma$-functions), an expansion by momentum regions,
or hypergeometric summation --- with high-precision numerical
evaluation used to corroborate, never to prove.
\end{definition}

\begin{example}[The phase-space volume]
\label{ex:volume}
The volume of the massless $n$-particle phase space in $d$
dimensions is fixed by Lorentz invariance and dimensional analysis
up to a constant, and iterating the two-body formula gives it in
closed form:
\begin{equation}
  \int \prod_{i=1}^{n} \dd^d p_i\, \delta^{(+)}(p_i^2)\;
  \delta^d\Big(q - \sum_i p_i\Big)
  \;=\;
  \big(q^2\big)^{\,(n-1)(d/2-1)-1}
  \left(\frac{\pi^{d/2-1}}{2}\right)^{\!n-1}
  \frac{\Gamma(d/2-1)^{\,n}}
       {\Gamma\big((n-1)(d/2-1)\big)\,
        \Gamma\big(n(d/2-1)\big)} ,
\end{equation}
a pure ratio of $\Gamma$-functions; other normalizations of the
measure differ by powers of $2\pi$ only.  (At $d = 4$ and $n = 3$
this reduces to the familiar $\Phi_3 = \pi^2 q^2/8$ in the
convention above.)  Three features make it the natural starting
point of the boundary analysis: it is the unit-power master of the
lowest sector; it requires no differential equation; and its
$\eps$-expansion supplies exact $\zeta$-values that propagate upward
through the inheritance conditions.  The genuinely new periods
appearing higher in the partial order are of the same character ---
cut integrals at an ordered limit, reducible to $\Gamma$-functions
or to low-dimensional Mellin--Barnes integrals.
\end{example}

\begin{example}[The two-dimensional system, concluded]
For the system of Example~\ref{ex:li2} the boundary vector has two
components per order.  At $w \to 0$ the residue matrix has
eigenvalues $\{0, -1\}$, offering local exponents $\{0, -\eps\}$; if
the underlying integral is finite and vanishing at $w = 0$, the mode
$w^{-\eps}$ is excluded and one constant is annihilated.  The other
is inherited: the first component is the volume-type subsector
master, already known.  The kernel of the constraint system is
trivial, no new period appears, and the solution
$\operatorname{Li}_2(w)$ carries no free constant --- in miniature,
exactly the accounting of the full calculation.
\end{example}

%======================================================================
\section{Behaviour at the edges of phase space}
\label{sec:endpoint}

The local analysis of Section~\ref{sec:local} identified, within the
global solution, the coefficient of each singular mode at each edge.
One further step connects these modes to the cross section.  The
cross section integrates the masters against smooth functions up to
the boundary, and near an endpoint the relevant integrals contain
modes of the form $(1-w)^{-1 + a\eps}$.  For such a mode the
expansion in $\eps$ and the integration against a test function do
not commute.  The correct object is the distributional identity
\begin{equation}
  (1-w)^{-1+a\eps} \;=\;
  \frac{1}{a\eps}\,\delta(1-w)
  \;+\;
  \sum_{k\ge0} \frac{(a\eps)^k}{k!}
  \left[\frac{\log^k(1-w)}{1-w}\right]_+ ,
  \label{eq:plus}
\end{equation}
where the \emph{plus distribution} is defined by
\begin{equation}
  \int_0^1 \dd w\; \big[g(w)\big]_+\, f(w)
  \;=\; \int_0^1 \dd w\; g(w)\,\big[f(w) - f(1)\big]
\end{equation}
for test functions $f$ smooth at $w = 1$.  The proof of
\eqref{eq:plus} is one line: integrate against $f$, add and
subtract $f(1)$,
\begin{equation}
  \int_0^1 \dd w\,(1-w)^{-1+a\eps} f(w)
  = \frac{f(1)}{a\eps}
  + \int_0^1 \dd w\,(1-w)^{-1+a\eps}\big[f(w) - f(1)\big] ,
\end{equation}
where the first term uses $\int_0^1 (1-w)^{-1+a\eps}\dd w =
1/(a\eps)$, convergent for $\mathrm{Re}(a\eps) > 0$ and continued
from there; the second integrand may be expanded in $\eps$ term by
term because $f(w) - f(1)$ vanishes at $w = 1$, making every
$\log^k(1-w)/(1-w)\cdot[f(w)-f(1)]$ integrable.  Modes
$(1-w)^{n+a\eps}$ with $n \le -2$ are treated by the same
subtraction carried to higher order and produce derivatives of
$\delta(1-w)$.

Had one expanded the \emph{mode} first, the left-hand side would
have become the series
$\frac{1}{1-w}\sum_k \frac{(a\eps)^k}{k!}\log^k(1-w)$, each term of
which is non-integrable at $w = 1$: the $\delta$-distribution ---
the very content the cross section needs --- is destroyed by
premature expansion.  For this reason the analytic evaluation
carries, next to the expanded polylogarithmic form of each master,
its exact local data: which exponents $n + a\eps$ occur at the
endpoint loci and with which coefficients.  The identity
\eqref{eq:plus} is applied to the unexpanded modes, and only
afterwards is everything expanded in the regulator.  Recovering the
endpoint behaviour is thus not an act of resummation but of
refraining: the residue matrices supplied the exponents, the
boundary analysis supplied the coefficients, and \eqref{eq:plus}
converts them into distributions.

%======================================================================
\section{Search and proof}
\label{sec:proof}

Nothing in the construction depends on how a transformation was
found.  An ansatz, a numerical fit, arithmetic modulo a prime ---
each produces a candidate and nothing more; what makes the candidate
a theorem is the identity
$T^{-1} A T - T^{-1}\dd T = \eps\sum_a R_a\,\dd\log\varphi_a$,
verified exactly.  The same asymmetry runs through the whole
subject: flatness establishes a connection, exact reconstruction
from constant residues establishes a canonical form, and a boundary
period enters only through a derivation in which numerics appear
afterwards, as a check on arithmetic, never as evidence.  Where a
verification is performed modulo primes at random points, its
failure probability is bounded and stated, and it is the stated
bound --- not an impression of reliability --- that the result
carries.

It is honest to record the state of the calculation these notes
describe.  At the time of writing, all but the three three-root
families carry canonical forms established in this way ---
eighty-eight of the ninety-one --- while the three-root families
await the Galois-algebraic treatment of
Section~\ref{sec:galois}; of the boundary periods, the lowest tier
is established and the remainder is an open campaign, as is the
endpoint mode analysis of the full inventory.  The method is
complete; the calculation, like every large calculation, is a
living object.

%======================================================================
\section{Notation and terminology}
\label{sec:summary}

\begin{description}[style=nextline]
\item[Integral family / representative family] A complete set of
inverse propagators \eqref{eq:familyint} / its representative
modulo cut-preserving redefinitions of the loop momenta.
\item[Sector, subsector] The subset of propagators present with
positive power; a subsector arises by deleting lines.
\item[Master integral] Basis element of a family's integrals modulo
integration-by-parts identities.
\item[Connection, flatness] The matrix one-form $A$ in
$\dd\vec I = A\,\vec I$; the integrability condition
$\dd A = A \wedge A$.
\item[Block, class] A strongly connected group of masters (a
diagonal block of the triangular connection); its equivalence class
under basis permutation composed optionally with
$v \leftrightarrow w$.
\item[Canonical form ($\eps$-form), letters, alphabet, residues]
The form $A = \eps\sum_a R_a\,\dd\log\varphi_a$ with constant
$R_a$; the functions $\varphi_a$; their set; the matrices $R_a$.
\item[Completion equation] The linear equation
\eqref{eq:completion} for the transformation of the entries below
the diagonal.
\item[Root content, rationalizing substitution] The rank of the
group of square classes in an alphabet; a rational change of
variables under which roots become rational.
\item[Galois algebra of the roots] The extension
$\Q(v,w,\eps)[r_1,\dots,r_k]/(r_i^2 - q_i)$ with automorphism group
$(\Z/2)^k$ acting by sign changes.
\item[Boundary period] A genuinely new constant surviving all
boundary conditions; the value of a cut phase-space integral at an
ordered limit.
\item[Local modes] The exponents $n + a\eps$ and logarithm powers
of the Frobenius solutions at a singular locus, read from the
residue matrix.
\end{description}

\paragraph{Further reading.}  Henn's paper introducing the
canonical form is short and repays being read first; Duhr's
lectures are the standard entry to iterated integrals and
polylogarithms, including the shuffle and stuffle structure these
notes pass over; the differential-equation method itself goes back
to Kotikov and to Gehrmann and Remiddi, and its algorithmic
reduction to canonical form to Lee (with Meyer's and Lee's
implementations the practical references); Anastasiou and Melnikov
introduced reverse unitarity for exactly the present kind of
phase-space problem; Besier, Wasser and Weinzierl treat the
rationalization of square roots systematically; Pak's algorithm is
the standard reference for the normal ordering of Symanzik
polynomials.

\end{document}
